\begin{lemma}[Common missing colors are closed after a PE]
Let the terminal step of a positive-rung ETT be a PE, let
$\varphi_n$ be its terminal coloring, let $R_n$ be the chosen closure of
$T_n(v_n)$, and let $T_n\vee R_n$ be their ordered union.  Assume
explicitly that $V(T_n)$ is elementary under the pre-step coloring
$\varphi_{n-1}$, as supplied by the lower-rung assertion (i).  Then
\[
 T_n\vee R_n\text{ is }
 \bigl(\overline{\varphi_n}(T_n)\cap
       \overline{\varphi_n}(R_n)\bigr)\text{-closed}.
\]
Equivalently, every color missing on both $T_n$ and $R_n$ has empty
colored boundary at their ordered union.
\end{lemma}
\begin{proof}
Put $X=V(T_n)$, $Y=V(R_n)$, and $U=V(T_n\vee R_n)$.
The PE certificate in \noderef{TashkinovSeries} and
\noderef{TashkinovSeriesStep} supplies the exact connecting edge
$f_n$, colors $\delta_n,\gamma_n$, and support $v_n$ at the last
positive slot.  The displayed pre-step elementarity hypothesis is
therefore exactly the conditional hypothesis of
\noderef{HkuParallelBoundaryStructure}.  That lemma gives two facts:
every $\alpha\in\overline{\varphi_n}(T_n)$ with
$\alpha\ne\delta_n$ has
$\partial_{\varphi_n,\alpha}(T_n)=\varnothing$, and every exact
$\delta_n$-boundary identity of $T_n$ is incident with a vertex strictly
preceding $v_n$.

By \noderef{HierarchyBase}, $R_n$ is a
\noderef{TashkinovClosure} of the ordered prefix $T_n(v_n)$.
Thus $Y$ is closed under $\varphi_n$.  Moreover, by
\noderef{TashkinovAugmentationSequence}, $R_n$ is obtained from
$T_n(v_n)$ through a finite, possibly empty, chain of steps from
\noderef{TashkinovAugmentation}.  At every such step the preceding tree
is a \noderef{TreeSequenceInitialSegment} of the succeeding tree.
Induction along the chain, including the empty-chain case, therefore makes
$T_n(v_n)$ an ordered initial segment of $R_n$.  Unpacking
\noderef{TreeSequenceInitialSegment} entry by entry and then using
\noderef{TreeSequenceVertexSet} gives
\[
             V(T_n(v_n))\subseteq Y.                              \tag{1}
\]
The parallel clause of \noderef{HierarchyBase}, together with
\noderef{OrderedTreeUnion}, identifies the hierarchy core with the exact
ordered union and gives
\[
             U=X\cup Y.                                           \tag{2}
\]

Fix
$\alpha\in\overline{\varphi_n}(T_n)\cap
\overline{\varphi_n}(R_n)$.  By \noderef{TreeMissingColors}, there is
a vertex $y\in Y$ at which $\alpha$ is missing.  Suppose that an exact
edge identity $g$ belongs to
$\partial_{\varphi_n,\alpha}(U)$.  By
\noderef{BoundaryColorEdges}, $g$ has color $\alpha$ and belongs to the
uncolored boundary of $U$.  By \noderef{BoundaryEdges}, its two ends may
be named $p$ and $q$ so that $p\in U$ and $q\notin U$.  Equation (2)
gives $p\in X\cup Y$.

First suppose $p\in Y$.  Since $Y\subseteq U$, we have $q\notin Y$.
Thus \noderef{BoundaryEdges} puts the same exact identity $g$ in the
boundary of $Y$, and its unchanged color $\alpha$ lets
\noderef{BoundaryColorEdges} put it in
$\partial_{\varphi_n,\alpha}(Y)$.  But $Y$ is closed, while $\alpha$
is missing at $y\in Y$; this is forbidden by
\noderef{ColoredVertexSetClosed}.  Hence $p\notin Y$, and consequently
$p\in X$ by (2).

Because $X\subseteq U$, the relation $q\notin U$ implies $q\notin X$.
The inside/outside characterization in \noderef{BoundaryEdges} therefore
puts $g$ in the boundary of $X$, and \noderef{BoundaryColorEdges}
upgrades this to $g\in\partial_{\varphi_n,\alpha}(X)$.

If $\alpha\ne\delta_n$, this contradicts the empty-boundary conclusion
of \noderef{HkuParallelBoundaryStructure}.  Suppose instead that
$\alpha=\delta_n$.  The other conclusion of that lemma supplies a tree
vertex $w\prec_{T_n}v_n$ incident with the exact identity $g$.
By the positional meaning of $\prec_{T_n}$ in
\noderef{TreeVertexOrder} and the definition of the prefix in
\noderef{TreeSequencePrefix}, the entry for $w$ occurs strictly before
the entry for $v_n$, so $w\in V(T_n(v_n))$.  Equation (1) gives
$w\in Y$.  The other endpoint of $g$ is $q$: indeed, $w\in X$ whereas
$q\notin X$, so the endpoint of $g$ lying in $X$ is the incident vertex
$w$.  Since $q\notin U$ also gives $q\notin Y$,
\noderef{BoundaryEdges} places $g$ in the boundary of $Y$, and
\noderef{BoundaryColorEdges} places it in
$\partial_{\varphi_n,\delta_n}(Y)$.  This again contradicts
\noderef{ColoredVertexSetClosed}, because $Y$ is closed and
$\delta_n=\alpha$ is missing at $y\in Y$.

We have ruled out every exact identity in
$\partial_{\varphi_n,\alpha}(U)$ for an arbitrary common missing color
$\alpha$.  Applying the definition \noderef{TreeColorSetClosed} to all
$\alpha\in\overline{\varphi_n}(T_n)\cap
\overline{\varphi_n}(R_n)$ proves the asserted color-set closedness.
\end{proof}
